\subsubsection{Geofencing}
\label{sec:geofencing}

Geofencing polega na wykrywaniu przekroczenia przez obiekt wirtualnej granicy
geograficznej wyłącznie w oparciu o dane lokalizacyjne. Wszystko, co zostało
opisane w poprzednich rozdziałach (komunikacja radiowa, integracja z Meshtastic,
wizualizacja na mapie), służy w ostateczności temu, żeby użytkownik dostał powiadomienie
w momencie, gdy węzeł przekroczy granicę wyznaczonego obszaru. W projekcie MeshTracker
granice definiowane są jako wielokąty o dowolnym kształcie, a monitorowanymi obiektami
są węzły sieci Meshtastic.

\paragraph{Model danych.}

Dane geofencingu przechowywane są w bazie Room za pomocą dwóch encji:
\texttt{Zone} i~\texttt{ZoneEvent} (tabele~\ref{tab:zone-entity}
i~\ref{tab:zoneevent-entity}).

\begin{table}[!h]
    \caption{Encja \texttt{Zone} - strefa geofencingu}
    \label{tab:zone-entity}
    \centering
    \begin{tabular}{|l|l|p{7cm}|}
        \hline
        \textbf{Pole} & \textbf{Typ} & \textbf{Opis} \\
        \hline
        \texttt{id}                 & \texttt{String (UUID)} & Klucz główny, generowany jako UUID v4 \\
        \texttt{name}               & \texttt{String}        & Nazwa nadana przez użytkownika \\
        \texttt{colorArgb}          & \texttt{Int}           & Kolor strefy (ARGB packed Int) \\
        \texttt{isActive}           & \texttt{Boolean}       & Czy strefa jest aktywnie monitorowana \\
        \texttt{verticesJson}       & \texttt{String}        & Wierzchołki wielokąta w formacie JSON \\
        \texttt{watchedNodeIdsJson} & \texttt{String}        & Lista ID węzłów do monitorowania w JSON \\
        \hline
    \end{tabular}
\end{table}

Wierzchołki wielokąta i lista monitorowanych węzłów przechowywane są jako serializowane
tablice JSON (np.\ \texttt{[\{"lat":52.1,"lon":21.0\},\ldots]}).

\begin{table}[!h]
    \caption{Encja \texttt{ZoneEvent} - zdarzenie wejścia/wyjścia ze strefy}
    \label{tab:zoneevent-entity}
    \centering
    \begin{tabular}{|l|l|p{6.5cm}|}
        \hline
        \textbf{Pole} & \textbf{Typ} & \textbf{Opis} \\
        \hline
        \texttt{id}               & \texttt{Long (autoincrement)} & Klucz główny \\
        \texttt{zoneId}           & \texttt{String}               & ID strefy  \\
        \texttt{nodeId}           & \texttt{String}               & ID węzła w chwili zdarzenia \\
        \texttt{eventType}        & \texttt{String}               & \texttt{"ENTER"} lub \texttt{"EXIT"} \\
        \texttt{timestampSeconds} & \texttt{Int}                  & Czas zdarzenia (Unix timestamp [s]) \\
        \hline
    \end{tabular}
\end{table}


\paragraph{Algorytm sprawdzania przynależności do strefy.}

Do wykrywania, czy punkt GPS leży wewnątrz wielokąta strefy, zastosowałem algorytm
\textbf{rzucania promienia} (ang.\ \textit{ray casting}) \footnote{\url{https://www.geeksforgeeks.org/cpp/point-in-polygon-in-cpp/}}.
Jest to jeden z~najszerzej stosowanych algorytmów rozwiązujących problem PIP (ang.\
\textit{Point-In-Polygon}) ze względu na prostotę implementacji i niską złożoność obliczeniową
$O(n)$, gdzie $n$ to liczba wierzchołków wielokąta~\cite{shimrat1962}. Uzasadnienie wyboru
akurat tej metody, oraz omówienie rozwiązań alternatywnych, zamieściłem na końcu tej sekcji.

Zasada działania jest następująca. Z badanego punktu $P = (x_P, y_P)$
wysyłany jest poziomy promień w kierunku rosnącej długości geograficznej $y \to +\infty$.
Zliczana jest liczba przecięć promienia z krawędziami wielokąta:

\begin{itemize}
    \item nieparzysta liczba przecięć $\Rightarrow$ punkt leży \textbf{wewnątrz} wielokąta,
    \item parzysta liczba przecięć $\Rightarrow$ punkt leży \textbf{na zewnątrz}.
\end{itemize}

Dla każdej krawędzi łączącej wierzchołki $X_i = (x_i, y_i)$
i~$Y_j = (x_j, y_j)$ sprawdzany jest warunek przecięcia z promieniem:
krawędź musi mieć jeden wierzchołek ściśle powyżej szerokości geograficznej punktu
$x_P$, a drugi poniżej lub równo z nią. Jeśli warunek jest spełniony, wyznaczana jest
długość geograficzna przecięcia:

\begin{equation}
    y_{\text{przecięcia}} = y_i +
        \frac{x_P - x_i}{x_j - x_i} \cdot (y_j - y_i)
    \label{eq:ray-cast-cross}
\end{equation}

Przecięcie jest liczone tylko jeśli $y_P < y_{\text{przecięcia}}$, tzn.\ leży
na prawo od badanego punktu (promień biegnie w prawo). Każde takie przecięcie neguje
bieżącą wartość flagi \texttt{inside}, która na początku wynosi \texttt{false}.


Listing~\ref{lst:geofence-contains} przedstawia pełną implementację algorytmu
w klasie \texttt{GeofenceChecker}.

\begin{lstlisting}[caption={Implementacja algorytmu ray casting w \texttt{GeofenceChecker.contains()}},
                   label={lst:geofence-contains},
                   basicstyle=\footnotesize\ttfamily,
                   xleftmargin=0pt]
fun contains(
    pointLat: Double,
    pointLon: Double,
    polygon: List<ZoneVertex>
): Boolean {
    if (polygon.size < 3) return false

    var inside = false
    var j = polygon.lastIndex

    for (i in polygon.indices) {
        val vi = polygon[i]
        val vj = polygon[j]

        // Czy krawedz przecina poziomy promien z punktu w prawo?
        val crossesRay =
            (vi.lat <= pointLat && vj.lat > pointLat) ||
            (vj.lat <= pointLat && vi.lat > pointLat)

        if (crossesRay) {
            // Dlugosc geograficzna przeciecia promienia z krawedzia
            val intersectLon = vi.lon +
                (pointLat-vi.lat)/(vj.lat-vi.lat)*(vj.lon-vi.lon)

            if (pointLon < intersectLon) {
                inside = !inside
            }
        }
        j = i
    }
    return inside
}
\end{lstlisting}

Implementacja obsługuje wielokąty \emph{wklęsłe} i~\emph{wypukłe} oraz~\emph{złożone} (z przecięciami własnymi). 
Nie jest jednak zoptymalizowana pod kątem przypadków brzegowych, takich jak punkt leżący dokładnie na krawędzi 
lub wierzchołku wielokąta - w takich sytuacjach wynik może być nieprzewidywalny. W praktyce, ze względu na 
niedokładność GPS, takie przypadki są pomijalnie rzadkie i nie stanowią problemu dla funkcjonalności aplikacji.

\paragraph{Wykrywanie zdarzeń wejścia i wyjścia.}

Sam algorytm odpowiada w sumie na pytanie \emph{,,czy węzeł jest teraz w strefie''}.
Aby wygenerować zdarzenia \texttt{ENTER} i \texttt{EXIT}, konieczne jest śledzenie
zmian stanu między kolejnymi aktualizacjami pozycji.


\texttt{ZoneMonitorService} utrzymuje słownik \texttt{previousNodeZoneMap}, 
mapujący~ID~węzła na zbiór ID stref, których węzeł znajdował się przy poprzedniej 
aktualizacji. Po każdym odebraniu broadcastu \texttt{ACTION\_NODE\_CHANGE} wykonywane są trzy kroki:

\begin{enumerate}
    \item Obliczany jest zbiór stref \texttt{currentZones} zawierających węzeł
          w~chwili obecnej (tylko spośród stref, które ten węzeł monitorują).
    \item Zdarzenia \texttt{ENTER} generowane są dla stref obecnych w \texttt{currentZones},
          ale nieobecnych w \texttt{previousZones}.
    \item Zdarzenia \texttt{EXIT} generowane są dla stref obecnych w \texttt{previousZones},
          ale nieobecnych w \texttt{currentZones}.
\end{enumerate}

Po przetworzeniu słownik jest aktualizowany w taki sposób, że jeżeli węzeł nie znajduje się w~żadnej strefie,
jego wpis jest usuwany, co minimalizuje zużycie pamięci.

\paragraph{\texttt{ZoneMonitorService} - monitorowanie w tle.}

Logika geofencingu działa w ramach \texttt{ZoneMonitorService} 
- usługi działające na pierwszym planie(ang.\ \textit{foreground service}) Androida. 
Jest to jedyna pewna metoda utrzymania procesu przy życiu w tle: system Android
agresywnie zamyka zwykłe procesy w tle, żeby oszczędzać baterię.
Usługi takie z widocznym powiadomieniem systemowym są chronione przed
zakończeniem~\cite{android-foreground-service}. Widoczne powiadomienie o działaniu serwisu
jest wymagane przez Androida i nie ma możliwości go ukryć - użytkownik zawsze widzi,
że aplikacja monitoruje strefy w tle.

Cykl życia serwisu jest następujący:
\begin{itemize}
    \item \textbf{Start}: \texttt{ZoneViewModel} woła \texttt{ZoneMonitorService.start()}
          gdy w bazie istnieje \textbf{co najmniej jedna aktywna strefa}.
    \item \textbf{Działanie}: serwis rejestruje \texttt{MeshtasticBroadcastReceiver}
          i nasłuchuje zmian stref z bazy Room przez \texttt{Flow}. Aktualizacja strefy
          w bazie jest propagowana do serwisu w czasie rzeczywistym bez potrzeby restartu.
    \item \textbf{Samozatrzymanie}: gdy \texttt{Flow} z Room zwróci pustą listę aktywnych
          stref (np.\ wszystkie zostały wyłączone lub usunięte), serwis wywołuje
          \texttt{stopSelf()} i kończy działanie.
    \item \textbf{Stop}: \texttt{ZoneViewModel} woła \texttt{ZoneMonitorService.stop()}
          gdy użytkownik \textbf{wyłącza ostatnią strefę}
\end{itemize}

\paragraph{Powiadomienia o przekroczeniu strefy.}

Wykrycie zdarzenia \texttt{ENTER} lub \texttt{EXIT} wyzwala wysłanie powiadomienia.
Powiadomienie zawiera nazwę węzła i nazwę strefy,wysyłane jest z priorytetem 
\texttt{PRIORITY\_HIGH} (kanał \texttt{CHANNEL\_ZONE\_ALERTS})
i~automatycznie znika po dotknięciu. 
Przy pierwszym uruchominiu aplikacji użytkownik musi jawnie przyznać uprawnienie \texttt{POST\_NOTIFICATIONS}
(obok lokalizacji i Bluetooth). Brak uprawnienia powoduje pominięcie powiadomienia bez
przerwania działania serwisu, co jest zachowaniem świadomym: serwis dalej wykrywa zdarzenia
i zapisuje je w bazie, jedynie nie wyświetla alertu.

\begin{figure}[!h]
    \centering
    \begin{subfigure}{0.48\linewidth}
        \centering
        \includegraphics[width=\linewidth]{uprawnieniaPowiadomienia.jpg}
        \caption{Prośba o uprawnienie - \texttt{POST\_NOTIFICATIONS}}
        \label{fig:uprawnienia-powiadomienia}
    \end{subfigure}
    \hfill
    \begin{subfigure}{0.48\linewidth}
        \centering
        \includegraphics[width=\linewidth]{strefaPowiadomienieWejscia.jpg}
        \caption{Powiadomienie o wejściu w strefę}
        \label{fig:strefa-powiadomienie-wejscia}
    \end{subfigure}
\end{figure}

Dotknięcie powiadomienia otwiera aplikację
MeshTracker na ekranie mapy.

\paragraph{Uzasadnienie wyboru i metody alternatywne.}

Zanim zdecydowałem się na ray casting, rozważałem kilka innych podejść, które warto
tu opisać - zarówno dlatego, że ich odrzucenie było nieoczywiste, jak i dlatego, że
w innych scenariuszach mogłyby być lepszym wyborem.

\textbf{Strefa okrągła (promień od centrum).}
Pierwszym pomysłem była strefa jako okrąg o środku $C = (x_C, y_C)$
i promieniu $r$. Sprawdzenie przynależności sprowadza się wtedy do prostego porównania:
\[
    d(P, C) \leq r \;\Leftrightarrow\; \text{punkt wewnątrz strefy}
\]
To rozwiązanie jest szybsze obliczeniowo i prostsze w implementacji - nie trzeba
przechowywać listy wierzchołków. Odrzuciłem je, bo działka na Mazurach ma kształt
zbliżony do wydłużonego wielokąta, a nie koła. Okrąg wystarczająco duży, żeby objąć
cały teren, wychodziłby daleko poza drogę przy granicy działki, przez co alarm byłby
spóźniony. Okrąg wystarczająco mały, żeby nie wychodzić za drogę, nie obejmowałby
części terenu. Wielokąt pozwala narysować granicę dokładnie wzdłuż ogrodzenia sąsiada,
polnej miedzy, albo linii lasu.

\textbf{Google Geofencing API.}
Wstępnie myślałem o skorzystaniu z gotowego 
geofencing API, które Android udostępnia właśnie do monitorowania przekroczeń granic
geograficznych. Szybko okazało się jednak, że to API monitoruje wyłącznie pozycję GPS
\emph{samego telefonu} i obsługuje tylko strefy okrągłe.


\textbf{Algorytm liczby owiniętej (ang.\ \textit{winding number}).}
Szukając literatury na temat PIP natknąłem się na algorytm winding number, który zlicza,
ile razy wielokąt owija się wokół badanego punktu. Jest numerycznie bardziej stabilny od
ray casting tzn. lepiej obsługuje przypadki brzegowe, gdzie punkt leży dokładnie na krawędzi
wielokąta. W moim przypadku to rozróżnienie jest jednak akademickie: GPS dokładny do kilku
metrów nigdy nie wyląduje na krawędzi wierzchołka zdefiniowanego przez użytkownika klikającego
na ekranie. Wybrałe Ray casting nad tym algorytmem, bo implementacja jest krótsza, prostsza do zrozumienia
i wystarczająca na potrzeby projektu. Złożoność obliczeniowa obu algorytmów wynosi $O(n)$.

\textbf{Indeksowanie przestrzenne (R-tree, quadtree).}
Gdybym chciał obsługiwać dziesiątki lub setki stref jednocześnie, liniowe sprawdzanie
każdej z nich przy każdym broadcastcie rzeczywiście stałoby się wąskim gardłem.
Struktury takie jak R-tree pozwalają szybko zawęzić kandydatów do stref, których obwiednia
prostokątna zawiera dany punkt. W MeshTrackerze spodziewam się co najwyżej kilku stref
i kilku węzłów - indeksowanie nie przyniosłoby żadnych mierzalnych korzyści, a znacznie
skomplikowałoby implementację.
